\section{Methods}
\label{sec:methods}

% MIDL 2027 CT&V call: "Submissions should clearly describe the study design, data
% provenance, reference standards, evaluation populations, statistical methods, and
% limitations." "Authors should distinguish internal testing from genuinely external
% testing" and "report performance across clinically relevant subgroups where appropriate."
% Coverage: design (3.1), provenance/reference standards/populations (3.2),
% internal vs external (3.1, 3.2), subgroups (3.2, 3.5), statistics (3.5), limitations (Sec. 5).
% Dataset provenance in 3.2 checked against the original dataset papers and Tudela 2024 (Nov 2026).

\subsection{Study design}
\label{sec:design}
This is a retrospective evaluation study. We propose no new segmentation model, but take trained models as they are and ask what geometric TTA changes in their outputs, at the level of averages, failures, detections and uncertainty.

We separate two kinds of testing. \emph{Internal testing} uses five-fold cross-validation on the training collection, with folds split by video (colonoscopy) or patient (MRI), so neither appears in both training and test data. \emph{External testing} applies the same checkpoints, unchanged, to three polyp collections and one polyp-free collection from other centres and devices. External data were never used for training, checkpoint selection or the choice of a model threshold, though the miss criteria were refined after seeing external results (Section~\ref{sec:eval}). Exactly one parameter was tuned, the low threshold $t$ of one recovery rule (Section~\ref{sec:tta}), on internal colonoscopy data only, so internal results for that rule are in-sample and external results are not. All data are public and de-identified, and code and software versions are in Appendix~\ref{app:repro}.

\subsection{Data, reference standards and evaluation populations}
\label{sec:data}

\paragraph{Internal colonoscopy.} CVC-ClinicDB \cite{bernal2015cvc} contains 612 frames recorded at Hospital Cl\'inic, Barcelona. The public release groups them into 29 sequences at $288\times384$ pixels. Each frame has a manually annotated polyp mask. % Source: polyp.grand-challenge.org/Databases ("612 still images from 29 different sequences")
Every frame contains a polyp. We assign whole sequences to five folds using the official frame-to-sequence table: for fold $k$, fold $k$ is the test set, fold $k+1$ (cyclically) the validation set and the rest training, so each sequence is a test unit exactly once and near-duplicate frames from one video never sit in both training and test data, which may be part of why our Dice is lower than values near 0.9 reported with random image-level splits \cite{fan2020pranet}. The release does not identify patients, and published descriptions report 31 sequences from 23 patients \cite{vazquez2017endoscene}, so our split prevents frame-level leakage but not every patient-level overlap.

\paragraph{Internal brain MRI.} BraTS 2021 \cite{baid2021brats} provides multi-institutional, multi-parametric MRI (T1, T1-Gd, T2, T2-FLAIR); we use the 1251 glioma patients of its training set. Tumour sub-regions were segmented by fusing automated methods, refined manually by neuroradiology experts and approved by board-certified neuroradiologists.
We use the enhancing tumour (ET, label 4) as the target; 1218 patients have ET and 33 do not. We split by patient into five folds in the same way. The ET-free patients have a glioma without enhancing tumour, and we count any predicted enhancing tumour in them as a false alarm.

\paragraph{External colonoscopy.} We use three polyp collections that were never seen in training: Kvasir-SEG \cite{jha2020kvasir} (1000 images), CVC-ColonDB \cite{bernal2012colondb} and ETIS-Larib \cite{silva2014etis}. Kvasir-SEG comes from Vestre Viken Health Trust, Norway; its polyp margins were outlined by an engineer and a medical doctor and checked by an experienced gastroenterologist. CVC-ColonDB \cite{bernal2012colondb} was recorded at Beaumont and St.\ Vincent's hospitals, Dublin \cite{tudela2024benchmark}, a different centre from our training collection, and ETIS-Larib comes from Lariboisi\`ere Hospital, Paris; both come with manual polyp masks. We take both from the test split distributed with PraNet \cite{fan2020pranet} (380 and 196 images), the split published polyp models report against; its CVC-ColonDB frame and sequence counts disagree with the original release by a route we could not establish (Appendix~\ref{app:repro}). ETIS-Larib's 196 images come from 34 sequences showing 44 polyps \cite{bernal2017challenge}, so several images can show the same polyp. We excluded CVC-300 and the PraNet copies of our training data and Kvasir-SEG. The shift ranges from mild (Kvasir-SEG) to severe (ETIS-Larib, which has many small polyps).

\paragraph{Polyp-free colonoscopy.} To measure false alarms we use the 1009 normal cecum frames of HyperKvasir \cite{borgli2020hyperkvasir} (class \emph{lower-gi-tract / anatomical-landmarks / cecum}), collected at B\ae{}rum Hospital, Vestre Viken Health Trust, Norway, and labelled by gastroenterologists or by junior doctors with expert review. The reference for every frame is empty. That label records an anatomical landmark rather than a read stating no polyp is present, so our false-alarm rates are upper bounds, used only to compare inference procedures on the same frames.

\paragraph{Subgroups.} Lesion size is a clinically relevant subgroup, so we split each external polyp collection into tertiles of reference polyp area in pixels, cut per collection. Pixel area confounds polyp size with camera distance and zoom, so these tertiles are not a size stratification in the clinical sense and no millimetre scale is recoverable from these releases. The polyp-free and ET-free populations above are the second subgroup.

These collections are distributed as images and most report no patient count: 23 patients for CVC-ClinicDB \cite{vazquez2017endoscene}, 13 for CVC-ColonDB \cite{tudela2024benchmark}, 44 polyps in 34 sequences for ETIS-Larib \cite{bernal2017challenge}, none for Kvasir-SEG or the cecum frames. The external evidence rests on well under 100 patients, and every count we report is an image count.

\subsection{Models and training}
\label{sec:models}
We trained three 2D architectures on CVC-ClinicDB, chosen to span a range of quality: a UNet \cite{ronneberger2015unet} with an ImageNet-pretrained ResNet-50 encoder \cite{he2016resnet}, a SegFormer with an ImageNet-pretrained MiT-B2 encoder \cite{xie2021segformer}, and a 2D SwinUNETR \cite{hatamizadeh2022swinunetr} trained from scratch. SwinUNETR from scratch on about 370 frames is weak by design: an example of a failed model, not a fair comparison of architectures. All 2D models saw only frames containing a polyp, so their false-alarm rates on normal frames are high in absolute terms and none is a deployment estimate; we use them only to compare inference procedures. Each architecture has 5 folds $\times$ 3 seeds, i.e.\ 15 checkpoints, all trained with binary cross-entropy plus Dice loss and AdamW on $352\times352$ inputs with geometric and brightness/contrast augmentation, keeping the checkpoint with the best validation Dice; the test fold was never used in training or selection (Appendix~\ref{app:training}).

For comparison with other ways of getting several predictions, we also trained a UNet with dropout ($p=0.2$) on encoder features and in the decoder, for Monte Carlo dropout \cite{gal2016dropout} (15 runs) and UNets on 50\% and 25\% of the training videos (5 folds at one seed each), giving a controlled series of weaker models.

On BraTS we trained a 3D SegResNet \cite{myronenko2019segresnet} (5 folds $\times$ 1 seed) on the four MRI sequences, with patch-based training and sliding-window inference (Appendix~\ref{app:training}). No model was retrained or fine-tuned after its training grid finished.

\subsection{Test-time augmentation and recovery rules}
\label{sec:tta}
\paragraph{Geometric TTA.} For each input we run the model on a set of transformed views, map each probability map back with the inverse transform, resample it to native resolution, average the maps and threshold the mean at 0.5. The 2D view sets are subgroups of the dihedral group $D_4$: TTA-2 is the identity and a horizontal flip, TTA-4 adds a 180$^\circ$ rotation and a vertical flip (the Klein four-group), TTA-8 is all of $D_4$; the 3D sets are the analogous axial subgroups with a depth flip, and every view was round-trip tested (Appendix~\ref{app:training}). The baseline is the identity view alone, and all methods in a comparison use the same forward passes of the same checkpoint.

\paragraph{Comparators.} A three-seed deep ensemble \cite{lakshminarayanan2017ensembles} averages the baseline maps of the three seeds of one fold, and the ensemble with TTA-4 averages all 12 maps. MC-Dropout averages 8 stochastic passes of the dropout UNet.

\paragraph{Recovery rules.} Averaging views can push a weak lesion below the 0.5 threshold even when some views fire on it. We test two rules that reuse the same views and need no extra passes. The \emph{any-view} rule declares a lesion present if the identity view or any TTA view is non-empty at 0.5. It then outputs the TTA mask if that is non-empty, else the identity-view mask, else the union of the non-empty view masks. It has no tuned parameter. The \emph{low-threshold} rule keeps the TTA mask when it is non-empty. When it is empty but the maximum of the mean map is at least $t$, it outputs the mean map thresholded at $t$. We report $t\in\{0.20, 0.30, 0.40\}$ and use $t=0.20$, fixed on internal UNet results before any lesion-free data were examined and then applied unchanged everywhere (Section~\ref{sec:design}). That procedure cannot in fact choose anything but the lowest $t$ on offer, because on polyp-only data filling an empty mask cannot lower Dice (Appendix~\ref{app:training}), so we state $t=0.20$ as a floor set by our grid and not as a selection, and report all three because the false-alarm cost rises steeply as $t$ falls. Neither rule ensures the output lies on the lesion. As a control we also apply the low-threshold rule to UNet's single view on the external sets.

\subsection{Evaluation and statistics}
\label{sec:eval}
\paragraph{Metrics.} We report Dice, the normalised surface distance (NSD) \cite{nikolov2021nsd}, computed on surface voxel counts, at a tolerance of 2\,px (2D) or 2\,mm (3D), and HD95 following the MONAI convention. Distances are computed at native resolution, in pixels for 2D and millimetres for BraTS, and because pixel sizes differ between colonoscopy collections we compare them across datasets only as relative changes. Distance metrics are defined only when both masks are non-empty, and paired comparisons use only cases non-empty under both methods, so an empty prediction cannot improve a distance. We count failures to find a lesion in three ways, from narrowest to broadest. A prediction is \emph{empty} if it has no foreground pixel. It has \emph{no overlap} if it shares no pixel with the reference (Dice $=0$); this adds non-empty predictions that lie entirely off the lesion. A lesion is \emph{missed} if the intersection over union (IoU) with the reference is below 0.10; this adds predictions that barely touch the lesion, and a lesion with IoU of at least 0.10 is \emph{localised}. We added the IoU criterion after seeing the first two counts on the external sets, so it is exploratory, and we report the whole sweep from 0 to 0.50 (Appendix~\ref{app:criteria}) so that no single bar carries the conclusion. An IoU of 0.10 is a bar for whether the mask is on the lesion, not for whether the segmentation is usable. On the polyp sets we count missed images, so a polyp shown in several images can count several times. A non-empty prediction on an empty reference is a \emph{false alarm}. A large failure is a case with HD95 $>50$\,px (CVC-ClinicDB) or $>10$\,mm (BraTS), or with no prediction from any seed. The 33 BraTS patients without enhancing tumour enter only the false-alarm analysis; every other BraTS analysis uses the 1218 with enhancing tumour.

\paragraph{Units and tests.} The unit of analysis is the video for CVC-ClinicDB (29), the patient for BraTS (1218 with ET) and the image for external sets, with the three seeds averaged per case internally so each video counts once. Externally, each image's metrics are averaged over an architecture's 15 checkpoints, and an image counts as empty, without overlap or missed if this holds for at least 8 of 15; the three votes are taken independently, so the nesting of the criteria holds per checkpoint but not per image. Internally the same counts use at least 2 of 3 seeds, while the HD95 analyses drop only frames with no prediction from any seed. This majority rule discards a spread larger than the effect it reports: on ETIS-Larib single checkpoints left between 1 and 29 images empty without TTA and between 5 and 42 with it, which is why the primary test below is per checkpoint. The 15 checkpoints are 5 folds $\times$ 3 seeds rather than independent replicates, so every result here describes this pool.

Metric comparisons use the paired Wilcoxon signed-rank test, with Holm correction across the metrics of each comparison (internal) or across metrics and TTA variants for each dataset and model (external). Confidence intervals are 95\% percentile bootstrap intervals resampling videos, patients or images. The counts that are this paper's subject are tested two ways. The primary test takes the checkpoint as the unit: for each of the 15 checkpoints we count the images meeting a criterion under each arm, and apply a two-sided sign test to the 15 paired differences with ties excluded, Holm-corrected across all 27 cells (3 models $\times$ 3 collections $\times$ 3 criteria). A deployed system is one checkpoint rather than a vote of 15, so this is the estimand the reader needs, and it does not depend on where the majority rule happens to cut. As a secondary, image-level view we apply the exact McNemar test to the majority votes, with one binary per image or patient per arm. Internal count comparisons carry Holm correction within each dataset. The external McNemar tests are \emph{not} corrected: Tables~\ref{tab:external} and~\ref{tab:criteria} hold 54 distinct such tests, since the 0 and 0.10 columns of Table~\ref{tab:criteria} repeat the no-overlap and missed columns of Table~\ref{tab:external}; the four sweep thresholds we do not tabulate add 36 more, and Section~\ref{sec:res-rules} adds 16. We therefore mark only $p<0.01$ there, read a single marked cell as a signal and not a finding, and do not read any $p$ above 0.01 as a signal at all; only a direction repeating across checkpoints, datasets and thresholds is a result. External images carry no sequence labels, so image-level tests treat them as independent, and we repeat the external comparisons on CVC-ColonDB and ETIS-Larib with a block bootstrap \cite{kunsch1989block} over contiguous blocks of consecutive image numbers; Kvasir-SEG carries no released sequence labels and was not blocked, so its intervals resample images. Those blocks are not the real sequences and are of equal length when real ones are not, so this is a sensitivity check and not a cluster-correct analysis. Clustering makes the image-level $p$-values optimistic, which bears on the rises and falls we do detect and not on the comparisons where we detect nothing: a non-rejection cannot be undone by widening an interval. Gains across size tertiles use the Kruskal--Wallis test, which shares the same independence assumption. Draw counts, Jeffreys intervals for external empty rates, the tail definitions and Oldham's correction \cite{oldham1962} for the baseline-error-versus-gain comparison are in Appendices~\ref{app:training} and~\ref{app:criteria}.

\paragraph{Uncertainty audit.} We score six per-case uncertainty signals computed from the same views, among them ROI entropy for the single view and for the TTA-4 mean map, TTA-4 variance in the ROI, and the fraction of pixels on which the binarised views do not all agree (full list in Appendix~\ref{app:training}); the ROI is the predicted foreground plus a 3-pixel band. Each signal is scored against the failures of its own prediction with AUROC and the area under the risk--coverage curve (AURC, risk $=1-$Dice), under three failure definitions: all failures, failures with a non-empty prediction (Dice $<0.5$), and empty predictions against clean cases, and signals are compared with a paired bootstrap.

\paragraph{Cost.} We estimate TTA latency as passes times the measured per-pass median at batch 1 on a laptop GPU. Batching the views would make it faster and the inverse transforms would make it slower, and we measured neither, so this is an estimate and not a bound in either direction (Appendix~\ref{app:deployment}).